% !PRINTED 255
\weberpage{255}

\begin{center}
\S.\,80.

\gesperrt{The Inertia of Forms of the Second Degree.}
\end{center}

By the theorem of the preceding section we are directed to the investigation of the homogeneous functions of the second degree with real coefficients. Concerning these functions we became acquainted in \S.\,58, under the name of the law of inertia, with a theorem which forms the foundation for what follows.

The theorem consisted in this: however one may transform a quadratic form of $m$ variables into a sum of positive and negative squares of linearly independent linear functions, which is possible in infinitely many different ways, the number of the positive squares, and likewise that of the negative ones, is always the same.

If we denote by $\pi$ the number of the positive squares and by $\nu$ the number of the negative ones, then $\pi + \nu$ is at most equal to $m$.

If we denote the difference $m - \pi - \nu$ by $\varrho$, then, when we represent the general quadratic form of $m$ variables as a sum of $m$ squares, $\varrho$ can be designated as the number of squares with vanishing coefficients.

The three numbers $\pi$, $\nu$, $\varrho$, of which none can be negative and whose sum is equal to the number of the variables,
\weq{1}{\pi + \nu + \varrho = m,}
are therefore unalterable for a given quadratic form.

We have to seek means of determining the numbers $\pi$, $\nu$, $\varrho$ from the coefficients of the quadratic form. If we apply these means to the Bezoutiant, we obtain criteria for inferring, from the coefficients of an equation, the number of its real roots.

Now, then, let, as in \S.\,56
\weq{2}{\varphi(x_1, x_2 \ldots x_m) = \sum_{1,m}^{i,k} a_{i,k} x_i x_k, \qquad a_{i,k} = a_{k,i}}
be a quadratic form of $m$ variables with \gesperrt{real} coefficients $a_{i,k}$.
